\chapter{Conception}

\section{Gestion du crédit, des débits et des écritures}

Dans cette section, « facturation » désigne uniquement la consommation du crédit d'écriture, sans paiement. Le sujet impose un crédit de 10\,000 caractères par rédacteur et par jour, commun à ses documents et à ses appareils, ainsi que l'impossibilité d'écrire au-delà de ce crédit et un blocage temporaire en cas d'acharnement \cite{sujet}. Il ne définit ni l'unité Unicode, ni l'heure de recharge, ni le détail des opérations facturables. Les règles qui suivent constituent donc un \emph{modèle proposé pour le MVP}, à confronter au fonctionnement réel de l'éditeur retenu et à valider avec le porteur du projet (Q-02, Q-04 à Q-07 et Q-09).

\subsection{Objet et unité du crédit}

Le crédit mesure le texte \emph{inséré et accepté} par un rédacteur pendant une journée. Il ne mesure ni les touches pressées, ni les échanges avec le système, ni la taille finale du document, ni les modifications reçues des autres rédacteurs. Le coût est cumulé pour un même utilisateur, quels que soient ses documents, ses sessions et ses appareils.

L'unité proposée est le \emph{point de code Unicode}, qui ne correspond pas toujours à un caractère perçu par le lecteur. Ainsi « Bonjour » contient sept points de code, « é » précomposé en contient un, tandis que « e » suivi d'un accent combinant en contient deux. Les espaces et les retours à la ligne saisis comptent ; les éléments ajoutés automatiquement pour le fonctionnement de l'éditeur ne comptent pas.

Le calcul proposé compte les points de code du texte effectivement inséré, sans transformer silencieusement sa forme Unicode. Deux suites visuellement identiques, comme les deux formes de « é » ci-dessus, peuvent donc avoir des coûts différents : ce choix doit être confirmé. Une règle commune doit aussi fixer comment compter les retours à la ligne, sans changer à l'insu des rédacteurs le texte ou les positions d'écriture.

\subsection{Coût proposé des opérations}

\begin{center}
\begin{tabularx}{\linewidth}{@{}L{43mm}Y@{}}
\toprule
Opération & Coût proposé \\
\midrule
Insertion ou collage & Nombre de points de code du nouveau texte inséré. \\
Suppression & Zéro, sans remboursement d'une insertion antérieure. \\
Remplacement & Nombre de points de code du texte nouvellement inséré. \\
Mise en forme seule, présence, déplacement du point d'insertion & Zéro. \\
Opération refusée ou retransmise après acceptation & Aucun débit définitif supplémentaire ; une retransmission doit retrouver le résultat de l'opération initiale. \\
Annulation & Zéro si elle supprime du texte ; coût des insertions si elle réinsère du texte. \\
Copie initiale d'une nouvelle version ou branche & Coût encore ouvert ; ne pas assimiler cette copie à une frappe du rédacteur. \\
Réinsertion volontaire d'un texte refusé & Coût des nouveaux points de code insérés, soumis aux droits et au crédit courants. \\
\bottomrule
\end{tabularx}
\end{center}

Ces coûts sont des propositions de conception, non des précisions énoncées par le sujet. Les remplacements et les annulations devront être vérifiés. Selon la décision du porteur, la restauration crée une branche indépendante : elle ne réécrit pas la version source. La gratuité de la copie initiale reste à valider ; une saisie ultérieure dans la nouvelle branche suit les règles ordinaires.

\subsection{Acceptation, plafond et refus}

Une frappe peut être préparée et affichée localement, puis reçue par le système, sans être encore acceptée. Le modèle distingue : \emph{préparée localement}, \emph{reçue}, \emph{admise} pour traitement, \emph{enregistrée} dans l'historique de l'éditeur, \emph{acceptée} et \emph{refusée}. L'admission ne vaut pas acceptation. Une confirmation d'acceptation ne doit être émise qu'après preuve que l'écriture durable et son débit sont cohérents. L'articulation exacte entre confirmation, retransmission et conservation des écritures reste à établir ; l'affichage local ne constitue pas cette preuve.

Soit $C(u,j)$ la consommation acceptée du rédacteur $u$ pour la journée $j$. L'invariant visé est :
\[
  C(u,j) \leq 10\,000.
\]
Une opération de coût $c$ ne peut être acceptée que si
$C(u,j)+c\leq 10\,000$.
Le refus entier est proposé lorsque l'opération dépasse le reliquat : avec 30 points de code disponibles, un collage de 50 points de code est refusé sans débit définitif. L'interface devra expliquer le refus et conserver le texte non accepté afin que le rédacteur puisse le récupérer. Une opération déjà acceptée, identifiée de manière stable, ne doit pas être débitée une seconde fois lors d'une retransmission.

\subsection{Journée, débit et abus}

La journée civile UTC est proposée : elle commence à 00\,h\,00 UTC et se termine juste avant le minuit suivant, sans report du reliquat. L'horloge du serveur fait foi et une opération appartient au jour de son acceptation, non à l'heure annoncée par le navigateur. Une fenêtre glissante de 24 heures depuis la première écriture produirait une autre règle ; ce n'est pas le choix proposé ici. La bascule de jour pendant une opération en cours devra être testée.

Le \emph{crédit} est une quantité cumulée par jour ; le \emph{débit} est un nombre de caractères ou d'opérations par unité de temps ; la \emph{protection contre les abus} constitue encore une autre règle. Le sujet n'impose pas d'étaler les 10\,000 caractères sur la journée : un collage volumineux peut être légitime. Aucun plafond arbitraire de caractères par seconde n'est fixé. Les seuils, la fenêtre d'observation, la durée et le périmètre d'un blocage temporaire restent à valider (Q-07).

\subsection{Identité, sessions et appareils}

Le crédit appartient à un utilisateur identifié de manière stable. Un rédacteur peut avoir plusieurs sessions simultanées ; fermer le navigateur, se reconnecter, changer de version ou créer une branche ne renouvelle pas son crédit. Le modèle conceptuel minimal comprend un \emph{rédacteur}, une \emph{session} avec expiration, une \emph{consommation journalière} et une \emph{opération d'écriture} identifiée de façon stable avec son état. Le crédit reste conservé par le système, indépendamment du navigateur. Un éventuel repère de navigateur ne prouve pas l'existence d'un appareil physique distinct.

Le sujet demande le nom et l'adresse e-mail sans compte classique \cite{sujet}. Le porteur ajoute une vérification par code envoyé à cette adresse avant d'établir l'identité métier stable \cite{decisions-porteur}. Plusieurs sessions peuvent alors partager le même crédit. La possession d'une adresse ne prouve pas qu'une personne ne dispose pas d'autres adresses. Il faudra vérifier qu'un même rédacteur conserve un crédit commun lors de plusieurs connexions, même si l'éditeur lui attribue plusieurs identifiants internes. Le choix de l'outil de vérification reste ouvert ; aucun mot de passe ni dispositif de reconnaissance du matériel n'est imposé.

\subsection{Cohérence du débit et décision à valider}

Un compteur atomique évite certains dépassements concurrents, mais ne rend pas à lui seul atomique l'ensemble « consommation du crédit + conservation de l'écriture ». Les échecs possibles sont : crédit débité sans écriture conservée, écriture acceptée sans débit, double débit après retransmission et dépassement entre appareils ou documents. Une réservation, un journal d'opérations sans double effet et une reprise après incident sont des pistes à étudier, sans garantie annoncée tant que le point de validation des écritures et les scénarios de panne n'ont pas été vérifiés.

Le refus d'une opération doit intervenir avant qu'elle ne soit diffusée comme acceptée et ne doit pas perturber les opérations suivantes. Le mécanisme de validation et l'ordre des confirmations seront examinés sur la version intégrée de l'éditeur. Un débit de crédit suivi séparément d'une sauvegarde ne démontre ni l'invariant de crédit ni la conservation des écritures acceptées. La qualification demandera des essais concurrents, des retransmissions et des pannes à chaque frontière d'acceptation.

\subsection{Diagramme de classes du crédit}
\umlfigure{credit-classes.png}{0.76}{Crédit d'écriture : objets conceptuels et frontière avec l'éditeur (cible).}{fig:credit-classes}
Ce diagramme sépare le compteur journalier, l'opération identifiée et la révision de l'éditeur. Le service de crédit et le journal d'opérations sont une cible de conception, pas des classes déjà opérationnelles.

\subsection{Diagramme d'activités du crédit}
\umlfigure{credit-activites.png}{0.72}{Contrôle proposé d'une écriture, du refus à la confirmation.}{fig:credit-activites}
Le parcours vérifie d'abord l'identité, le droit et la version ; il calcule ensuite le coût des seuls points de code insérés. Le débit durable, le traitement de l'échec et la confirmation demandent encore un protocole démontré.

\subsection{Diagramme de séquence du crédit}
\umlfigure{credit-sequence.png}{0.78}{Échanges visés pour écrire avec crédit suffisant ou refuser.}{fig:credit-sequence}
La frontière entre le stockage métier et les révisions de l'éditeur est volontairement visible. La séquence indique un besoin de réconciliation ; elle ne prétend pas qu'une transaction atomique existe déjà entre les deux systèmes.

\needspace{15\baselineskip}
\section{Historique, versionnage et branches}\label{sec:conception-historique}

Les règles de cette section précisent l'usage du journal décidé par le porteur du projet \cite{decisions-porteur}. Elles complètent le sujet initial. Le code réalise les parcours locaux de branches, de versions et de clôture ; les garanties de panne et de charge restent à démontrer.

\subsection{Termes et organisation des documents}

Une \emph{révision technique} est un état enregistré au fil des modifications de l'éditeur. Une \emph{version métier} (V1, V2, etc.) est une étape de travail explicitement ouverte, puis figée. Elle peut contenir de nombreuses révisions techniques : chaque frappe ne crée pas une nouvelle version métier. Une \emph{branche} rassemble une succession de versions métier et possède son propre historique. Une version appartient à une seule branche ; une branche ne possède au plus qu'une version ouverte à l'édition. Pendant une clôture, sa version active est en état FERMETURE et la suivante n'est pas encore éditable.

Le projet d'édition porte le propriétaire unique, le fuseau IANA et la prochaine échéance. Chaque branche porte ses droits et sa version active ; chaque version porte son numéro, son état, sa référence de pad et sa révision finale lorsqu'elle est figée. Une provenance désigne la branche, la version et la révision sources. Une clôture possède une cause, une échéance et un avancement. Ces objets existent dans le schéma métier actuel, sous des noms anglais ; la table historique 	exttt{documents} joue le rôle de projet d'édition. Elle sera renommée seulement par une migration préservant les données.

Par exemple, la branche A contient V1 figée puis V2 éditable. Si son propriétaire restaure une révision de V1, le système crée la branche B et sa première version éditable à partir de ce texte. La provenance de B mentionne A, V1 et la révision choisie. V1 et V2 de A ne changent pas ; les modifications ultérieures de A et de B ne se propagent pas entre les branches. Même avec un texte initial identique, les deux documents de travail sont distincts. Le MVP ne prévoit pas de fusion automatique.

\subsection{Autorisations et édition en ligne}

\begin{tabularx}{\linewidth}{@{}L{28mm}@{\hspace{2mm}}L{36mm}@{\hspace{2mm}}L{35mm}@{\hspace{2mm}}Y@{}}
\toprule
Rôle sur la branche & Lire ses versions et révisions & Écrire dans la version active & Clore, restaurer et attribuer les droits \\
\midrule
Propriétaire & Oui. & Oui, tant qu'elle est OUVERTE. & Oui ; la restauration crée une nouvelle branche. \\
\addlinespace
Collaborateur en écriture & Oui. & Oui, tant qu'elle est OUVERTE et que le droit subsiste. & Non. \\
\addlinespace
Collaborateur en lecture & Oui. & Non. & Non. \\
\addlinespace
Personne sans droit & Non ; un lien public de PDF suit une règle distincte. & Non. & Non. \\
\bottomrule
\end{tabularx}

Le droit de lecture porte sur la branche et son historique. Le droit d'écriture vise uniquement la version active ouverte ; une version figée reste en lecture seule. La bascule vers V2 conserve les droits de la branche, y compris leurs révocations, sans nouvelle invitation. Une restauration n'hérite pas automatiquement des droits de la source : le propriétaire attribue les accès à la branche créée. Le serveur doit vérifier ces règles avant d'admettre une écriture ; un verrouillage visuel seul ne suffit pas.

La saisie suppose une connexion opérationnelle au serveur d'édition. Lorsqu'une coupure est détectée, l'interface doit cesser d'accepter de nouvelles saisies. La détection n'est pas instantanée : le texte affiché localement pendant cet intervalle n'est pas pour autant accepté. À la reconnexion, le système doit vérifier de nouveau la version active, les droits et le crédit avant toute écriture ; il ne doit pas rejouer automatiquement une modification devenue interdite.

\subsection{Fermeture manuelle et passage quotidien}

\umlfigure{cycle-version-metier.png}{0.72}{Cycle métier d'une version ; la suivante est un autre document de travail.}{fig:cycle-version-metier}

Lors d'un gel manuel, le serveur enregistre la demande et fixe une échéance trois secondes plus tard selon son horloge. Il informe les clients, qui désactivent les nouvelles saisies. La marge permet l'admission d'opérations reçues par le serveur avant l'échéance, sous réserve des droits et du crédit ; une heure déclarée par le navigateur ne fait pas foi. Après l'échéance, toute opération encore non admise est refusée explicitement. Les trois secondes bornent la fenêtre d'admission, pas le temps total de persistance ou de création de V2.

La clôture suit un ordre : fermer l'admission, résoudre chacune des opérations admises, fixer la révision finale, figer V1, créer le nouveau document de travail depuis ce contenu, puis activer V2. Une opération admise peut encore être refusée avec une réponse explicite ; une opération déjà acceptée ne peut pas être abandonnée. V2 ne devient éditable qu'une fois la source finale déterminée. La reprise après incident doit reconnaître la clôture par son identifiant et ne créer qu'un seul successeur actif.

Le même mécanisme se déclenche à la fin de la journée civile du propriétaire, même s'il est absent. Son fuseau IANA est conservé explicitement ; le backend calcule et recherche périodiquement les échéances, sans minuteur permanent par document. Si la fermeture commence à minuit local, l'admission cesse trois secondes plus tard, puis la finalisation peut prendre davantage de temps : V1 n'est pas un instantané strict de minuit. Les droits de branche demeurent applicables à V2 lorsqu'elle est prête. Cette journée de versionnage, la journée UTC proposée pour le quota du rédacteur et la durée de ses sessions sont trois notions distinctes. La clôture ne recharge jamais le crédit.

\subsection{Écritures refusées, invariants et décisions restantes}

Une opération trop tardive est refusée avec une explication. Une proposition simple pour le MVP est un panneau ou un export du texte local, afin que le rédacteur puisse récupérer ce qu'il souhaite réinsérer lui-même dans V2. Cette réinsertion passe à nouveau par les droits, le crédit et l'acceptation habituelle. Le texte local complet peut déjà contenir des contributions acceptées : il ne faut pas le copier intégralement dans V2 comme s'il représentait seulement le refus. Le client local conserve maintenant cette copie dans le navigateur après rechargement ; les essais de reconnexion longue restent à faire.

Les invariants visés sont : au plus une version OUVERTE par branche ; aucune écriture ordinaire dans une version FIGÉE ; aucune perte d'une écriture acceptée pendant la clôture ; aucune nouvelle admission après l'échéance ; aucune activation du successeur tant que la révision finale reste indéterminée ; aucune propagation automatique entre branches ; aucune recharge de crédit par changement de version ; un seul successeur pour une même clôture, y compris après reprise.

Restent à valider la rétention des révisions, la création quotidienne d'une nouvelle version lorsque le texte n'a pas changé, la facturation de la copie initiale et la façon de distinguer précisément les opérations déjà en transit lorsque le gel commence. La preuve des garanties d'admission et de persistance dépendra de l'intégration réelle avec l'éditeur.


\subsection{Diagramme de classes des versions}
\umlfigure{version-classes.png}{0.78}{Projet, branches, versions, provenance et clôture.}{fig:version-classes}
La propriété est portée par le projet. Une branche ne possède pas un propriétaire autonome ; elle conserve ses permissions au passage à V2. Chaque version métier référence un pad distinct et conserve ses révisions techniques dans ce pad.

\subsection{Diagramme d'activités de la clôture}
\umlfigure{version-activites.png}{0.72}{Clôture ordonnée avec reprise après incident.}{fig:version-activites}
Les trois secondes bornent l'admission. La barrière fixe la révision finale avant la création du successeur ; une panne laisse une opération à reprendre et ne justifie pas l'activation prématurée de V2.

\subsection{Diagramme de séquence de la clôture}
\umlfigure{version-sequence.png}{0.78}{Échanges entre propriétaire ou horloge, backend, éditeur et clients.}{fig:version-sequence}
La restauration suit une autre demande propriétaire : elle copie une révision choisie dans une nouvelle branche du même projet. Les changements des deux branches restent indépendants et les collaborateurs ne sont pas invités implicitement.

\section{Création et partage des projets d'édition}\label{sec:conception-partage}

\subsection{Règles et responsabilités}
Le terme \emph{projet d'édition} désigne l'ensemble appartenant à un seul propriétaire ; l'ancienne table nommée \texttt{documents} représente actuellement cet objet et ne crée pas une seconde entité concurrente. La création idempotente produit la branche initiale, V1 ouverte et un pad indépendant. Le propriétaire contrôle toutes les branches. Un collaborateur reçoit READ ou WRITE uniquement sur la branche invitée ; WRITE inclut READ. Aucun droit n'autorise l'écriture dans une version figée.

Le propriétaire enregistre une invitation pour une adresse, un droit et une branche. Le lien contient un jeton aléatoire dont seule l'empreinte figure dans l'invitation persistée. Le lien n'accorde rien : le destinataire doit se connecter, vérifier la même adresse par code, puis accepter explicitement. L'acceptation crée ou modifie la permission de branche. Les états de l'invitation (PENDING, ACCEPTED, REVOKED, EXPIRED) sont distincts de ceux de l'envoi (QUEUED, SENT, FAILED). Un renvoi change le jeton et invalide l'ancien lien. Une invitation acceptée ne peut pas restaurer un droit révoqué par son rejeu.

\subsection{Diagramme de classes du partage}
\umlfigure{partage-classes.png}{0.78}{Projet propriétaire, permission de branche et invitation distincte.}{fig:partage-classes}
Les noms français désignent des concepts métier ; le code existant utilise notamment \texttt{BusinessService} et \texttt{InvitationService}. La relation entre utilisateur et projet porte la propriété unique. Une invitation en attente est séparée de la permission effectivement accordée.

\subsection{Diagramme d'activités de l'invitation}
\umlfigure{partage-activites.png}{0.72}{De la création de l'invitation à l'acceptation ou au refus.}{fig:partage-activites}
L'activité montre les refus dus à la mauvaise adresse, à l'expiration et à la révocation. Un échec d'envoi laisse l'état FAILED ; le propriétaire peut demander un renvoi qui renouvelle le lien. Le mode local garde les messages dans une boîte de test ; l'envoi SMTP externe n'est pas une livraison exactement une fois.

\subsection{Diagramme de séquence du partage}
\umlfigure{partage-sequence.png}{0.78}{Invitation d'une adresse encore inconnue et création du droit après vérification.}{fig:partage-sequence}
La création d'un compte avant l'invitation n'est pas nécessaire. Le service d'authentification prouve l'adresse ; le backend métier demeure responsable de la permission. L'acceptation sous verrou et l'unicité de la permission empêchent les doublons lors de requêtes concurrentes.

\subsection{Invariants et limites}
Une adresse invitée n'obtient aucune métadonnée d'une autre branche. Le propriétaire peut modifier ou révoquer les collaborateurs ; un collaborateur ne peut pas inviter. Les droits de branche survivent à V2 mais ne sont jamais copiés sur une branche restaurée. Le relais SMTP externe, les messages perdus entre enregistrement et envoi ainsi que les reprises après panne réclament encore une qualification. L'authentification déclarative historique reste uniquement un adaptateur de développement ; l'acceptation d'une invitation exige une session issue d'un code e-mail vérifié.
